\documentclass[10pt,a4paper]{amsart}
\usepackage[margin=19mm]{geometry}
\usepackage{amsmath,amssymb,mathtools}
\usepackage[scr=rsfs,cal=euler]{mathalfa}
\usepackage{xfrac}
\usepackage[unicode,hidelinks,bookmarks=false]{hyperref}
\newcommand{\ie}{i.e.\ }
\newcommand{\cf}{cf.\ }
\newcommand{\resp}{respectively\ }
\newcommand{\Subsection}{\subsection}
\providecommand{\nobreakdash}{\nobreak\hbox{-}}
\setlength{\emergencystretch}{2em}
\multlinegap0pt
\usepackage{amssymb}
\usepackage{mathtools}
\usepackage[shortlabels]{enumitem}
\newtheorem{theorem}{Theorem}
\newcommand{\C}{\mathbb{C}}
\newcommand{\prof}{\operatorname{prof}}
\title{Pl\"ucker coordinates of finite-dimensional subspaces of $\ell^p$ and its direct sums:\\ summability, reconstruction, stratification}
\author{David Victor Feldman}
\date{Source: arXiv:2608.00983v1 (CC BY 4.0)}
\begin{document}
\maketitle

\noindent\emph{Source and licence.} Pl\"ucker coordinates of finite-dimensional subspaces of $\ell^p$ and its direct sums:\\ summability, reconstruction, stratification. Version arXiv:2608.00983v1 (CC BY 4.0), distributed by arXiv under the Creative Commons Attribution 4.0 International licence, http://creativecommons.org/licenses/by/4.0/, as recorded in the arXiv licence metadata for this exact version and shown on the abstract page https://arxiv.org/abs/2608.00983v1. Full licence notice: \texttt{licenses/arxiv-2608.00983-CC-BY-4.0.txt}.\par\smallskip

\noindent\textit{Editorial scope.} Counted unit: the article's theorem thm:interval (source0.tex lines 276--278). The article's complete original proof of it is delivered with it (source0.tex lines 280--288, one \texttt{proof} environment of 9 lines). No mathematics is written, repaired or re-derived: the statement and the proof are the article's own lines, in the article's own notation, and no retained line is substituted or re-flowed.  Retained uncounted as the source-local context the proof needs: lines 142--144.  Nothing was omitted from any retained span, so the package carries no omission record.  No retained line contains a citation, so the wrapper carries no bibliography and no bibliography entry.  No retained line contains a figure environment, an image-inclusion command or drawing code, so no image is delivered and the package declares no asset dependency.  Editorial formatting only, in addition: an AMS \texttt{amsart} preamble that restates the article's own declarations and macros used by the retained lines, namely the environments \texttt{theorem} on separate counters, together with the standard packages those lines need.  The retained environments print in this document as Theorem 1 in order of appearance, whereas the article prints the same items with the numbering of its own full document.  The complete native source of the article as distributed in the arXiv e-print for this exact version is bundled unchanged as \texttt{sources/206584/source0.tex}.
\begin{equation}\label{eq:plucker}
\sum_{t=1}^{n+1}(-1)^t\,\tilde p_{\,R\,s_t}\;\tilde p_{\,s_1\cdots \widehat{s_t}\cdots s_{n+1}}\;=\;0 .
\end{equation}

\begin{theorem}[Interval theorem]\label{thm:interval}
$\prof(V)=\{(a,n-a):\ \alpha\le a\le n-\gamma\}$: an interval with no gaps, determined by the two integers $(\alpha,\gamma)$.
\end{theorem}

\begin{proof}
Represent $V$ by $M=[M_A\mid M_B]$; block $(a,b)$ is nonzero iff some nonvanishing $n\times n$ minor uses exactly $a$ columns from $A$.

\emph{Exchange, from the relations.} Let $p_I\ne0\ne p_J$ and $i\in I\setminus J$. Instantiate \eqref{eq:plucker} with $R=I\setminus\{i\}$ and $S=J\cup\{i\}$: the terms are $\pm\,p_I\,p_J$ and, for $j\in J$, $\pm\,p_{(I\setminus i)\cup j}\ p_{(J\setminus j)\cup i}$. Were every exchanged minor $p_{(I\setminus i)\cup j}$ to vanish, the relation would force $p_Ip_J=0$; so some exchange survives, and any surviving $j$ lies in $J\setminus I$ --- a value already in $I\setminus\{i\}$ would repeat a column, and $j=i$ is excluded by $i\notin J$. Base exchange for the column matroid is thus a consequence of the quadratic relations themselves; no matroid theory is imported.

\emph{Discrete intermediate value.} Iterate: each exchange replaces an element of $I\setminus J$ by an element of $J\setminus I$, strictly shrinking $I\setminus J$, so the process connects $I$ to $J$ through nonvanishing minors while changing the $A$-column count by at most one per step. Every count between the two is achieved.

\emph{Endpoints.} The largest achievable $A$-count is the rank of the $A$-columns: at most, because the $A$-columns of a nonvanishing minor are linearly independent; at least, because a maximal independent set of $A$-columns extends within the columns --- which span $\C^n$, some minor being nonvanishing --- to an independent $n$-set (Steinitz exchange), and $n$ independent columns form a nonvanishing minor. Since $\ker(\pi_A|_V)=V\cap\ell^qB$, the $A$-column rank is $n-\gamma$; symmetrically for $B$, the least achievable $A$-count is $n-(n-\alpha)=\alpha$. Hence $\prof(V)=\{(a,n-a):\ \alpha\le a\le n-\gamma\}$.
\end{proof}

\end{document}
